% Discrete Mathematics (MAL2190), IIT Jodhpur
% Common project report template: Build a Game, Discover the Mathematics
% Upload this file to Overleaf and compile with pdfLaTeX (twice).
% This source is self-contained. An optional IITJ_COLOURED.png may be placed
% in the same folder to display the institute logo.
%
% STUDENT INSTRUCTIONS
% 1. Edit the four metadata commands below and all bracketed placeholders.
% 2. Keep the section order and common formatting. Add pages as needed.
% 3. Delete unused example rows and add your own figures/tables where useful.
% 4. Set \showguidancefalse before final submission to hide guidance boxes.
%    This does NOT hide unfilled placeholders: replace every one of them.
% 5. The page breaks below organise the preview; move/remove them as needed.
% 6. No page limit, deadline, or marking scheme is specified by this template.

\documentclass[11pt,a4paper]{article}
\usepackage[T1]{fontenc}
\usepackage{lmodern}
\usepackage[a4paper,margin=1in,headheight=15pt]{geometry}
\usepackage{amsmath,amssymb,amsthm}
\usepackage{graphicx}
\usepackage[table]{xcolor}
\usepackage{array,tabularx,longtable,booktabs}
\usepackage{enumitem}
\usepackage{fancyhdr,lastpage}
\usepackage{titlesec}
\usepackage[hidelinks]{hyperref}

% ---------- EDIT THESE DETAILS ----------
\newcommand{\GameTitle}{[Title of the game or puzzle]}
\newcommand{\GroupNumber}{[Group number]}
\newcommand{\SubmissionDate}{[DD Month YYYY]}
\newcommand{\MathConcepts}{[Main discrete mathematics concept(s)]}
% ----------------------------------------

\newif\ifshowguidance
\showguidancetrue % Change to \showguidancefalse for the completed report.
\definecolor{guidetext}{gray}{0.32}
\definecolor{guidebg}{gray}{0.95}
\newcommand{\placeholder}[1]{\textit{[#1]}}
\newcommand{\guidance}[1]{%
  \ifshowguidance
    \par\smallskip\noindent
    \begingroup\setlength{\fboxsep}{8pt}%
    \colorbox{guidebg}{\parbox{\dimexpr\linewidth-2\fboxsep\relax}{%
      \small\color{guidetext}#1}}\endgroup\par\smallskip
  \fi}
\newcolumntype{Y}{>{\raggedright\arraybackslash}X}
\newcolumntype{P}[1]{>{\raggedright\arraybackslash}p{#1}}
\setlength{\parindent}{0pt}
\setlength{\parskip}{5pt}
\setlength{\emergencystretch}{2em}
\renewcommand{\arraystretch}{1.3}
\setlist{leftmargin=*,itemsep=3pt,topsep=4pt,parsep=0pt}
\titleformat{\section}{\large\bfseries}{\thesection.}{0.5em}{}
\titleformat{\subsection}{\normalsize\bfseries}{\thesubsection}{0.5em}{}
\titlespacing*{\section}{0pt}{13pt}{5pt}
\titlespacing*{\subsection}{0pt}{9pt}{3pt}
\raggedbottom
\pagestyle{fancy}
\fancyhf{}
\fancyhead[L]{\small Discrete Mathematics (MAL2190)}
\fancyhead[R]{\small Project Report}
\fancyfoot[L]{\footnotesize IIT Jodhpur}
\fancyfoot[R]{\footnotesize Page \thepage\ of \pageref{LastPage}}
\renewcommand{\headrulewidth}{0.4pt}
\renewcommand{\footrulewidth}{0.4pt}
\hypersetup{
  pdftitle={Discrete Mathematics - Project Report Template},
  pdfsubject={Build a Game, Discover the Mathematics}
}

\begin{document}
\thispagestyle{plain}

\begin{center}
  \IfFileExists{IITJ_COLOURED.png}{%
    \includegraphics[width=1.6cm]{IITJ_COLOURED.png}\par\vspace{6pt}}{}
  {\large\bfseries Indian Institute of Technology Jodhpur\par}
  \vspace{4pt}
  {\large\bfseries Discrete Mathematics (MAL2190)\par}
  \vspace{10pt}
  \rule{\textwidth}{0.6pt}\par
  \vspace{6pt}
  {\LARGE\bfseries Project Report\par}
  \vspace{5pt}
  {\large Build a Game, Discover the Mathematics\par}
  \vspace{6pt}
  \rule{\textwidth}{0.6pt}\par
  \vspace{13pt}
  {\Large\bfseries \GameTitle\par}
\end{center}

\vspace{5pt}
\begin{tabularx}{\textwidth}{@{}P{3.8cm}Y@{}}
\textbf{Group number} & \GroupNumber \\
\textbf{Submission date} & \SubmissionDate \\
\textbf{Concept(s)} & \MathConcepts \\
\end{tabularx}

\vspace{9pt}
\textbf{Group members (10 students)}

\begin{tabularx}{\textwidth}{@{}P{1cm}Y P{4cm}@{}}
\toprule
\textbf{No.} & \textbf{Name} & \textbf{Roll number} \\
\midrule
1 & \placeholder{Name of member 1} & \placeholder{Roll number} \\
2 & \placeholder{Name of member 2} & \placeholder{Roll number} \\
3 & \placeholder{Name of member 3} & \placeholder{Roll number} \\
4 & \placeholder{Name of member 4} & \placeholder{Roll number} \\
5 & \placeholder{Name of member 5} & \placeholder{Roll number} \\
6 & \placeholder{Name of member 6} & \placeholder{Roll number} \\
7 & \placeholder{Name of member 7} & \placeholder{Roll number} \\
8 & \placeholder{Name of member 8} & \placeholder{Roll number} \\
9 & \placeholder{Name of member 9} & \placeholder{Roll number} \\
10 & \placeholder{Name of member 10} & \placeholder{Roll number} \\
\bottomrule
\end{tabularx}

\vspace{10pt}
\guidance{\textbf{Using this template.} Replace all bracketed text with your
group's work. Keep the main section headings and formatting so the reports can
be combined for the course website. Add space and pages where needed; the
preview length is not a prescribed page limit. Remove guidance boxes before
submission. Explain the game so that a reader can recreate and play it.}

\newpage
\section{The Game or Puzzle}

\subsection{Overview and objective}
\placeholder{Introduce the game in a short paragraph. State what the players
are trying to achieve and what makes the activity interesting.}

\subsection{Players, materials, and initial setup}
\begin{itemize}
  \item \textbf{Number of players:} \placeholder{Individual, two-player, team,
  or other format.}
  \item \textbf{Materials or platform:} \placeholder{List the board, cards,
  counters, software, or other items needed. For a digital game, explain how
  to access and start it.}
  \item \textbf{Initial setup:} \placeholder{Describe the starting arrangement
  and how the first player or first move is chosen, if applicable.}
\end{itemize}

\subsection{Rules and how to play}
\guidance{Write the rules in the order a new player needs them. State the legal
moves, restrictions, turn order, and ending condition. Define winning, losing,
scoring, ties, or puzzle completion wherever relevant.}
\begin{enumerate}
  \item \placeholder{First rule or step.}
  \item \placeholder{Second rule or step.}
  \item \placeholder{Further rules and restrictions; add items as needed.}
  \item \placeholder{When the game ends and how the outcome is decided.}
\end{enumerate}

\subsection{Illustrative play-through}
\placeholder{Show one short example from a valid starting state through the
main moves to an outcome. Use a labelled diagram, photograph, or table if it
helps. Explain why each move is allowed.}

% Example figure syntax, if needed:
% \begin{figure}[htbp]
%   \centering
%   \includegraphics[width=0.7\linewidth]{game_setup.png}
%   \caption{Initial setup of the game.}
%   \label{fig:setup}
% \end{figure}

\section{The Mathematical Question}
\placeholder{State the precise mathematical question that arises from the
game. What do you want to determine, count, characterise, or prove?}

\textbf{Connection to discrete mathematics.}
\placeholder{Name the relevant concept(s) and explain how the game represents
them. Define the symbols, objects, and assumptions used in your analysis.}

\newpage
\section{Experimentation and What We Discovered}

\subsection{Method and observations}
\placeholder{Explain how you explored the game: which starting positions,
sizes, strategies, or rule variations you tested, and how you recorded the
outcomes. If you used code, describe what it checked.}

\guidance{Record actual trials. Distinguish one play-through from an exhaustive
check of all possibilities. Add or modify rows and columns to suit your game.}

\begin{tabularx}{\textwidth}{@{}P{1cm}Y Y Y@{}}
\toprule
\textbf{Trial} & \textbf{Starting case} & \textbf{Moves or strategy tested}
  & \textbf{Outcome / observation} \\
\midrule
1 & \placeholder{Case} & \placeholder{What you tried} & \placeholder{Result} \\
2 & \placeholder{Case} & \placeholder{What you tried} & \placeholder{Result} \\
3 & \placeholder{Case} & \placeholder{What you tried} & \placeholder{Result} \\
\bottomrule
\end{tabularx}

\subsection{Patterns and discoveries}
\placeholder{Describe the patterns you noticed and the evidence for them.
Mention any surprising outcomes, exceptions, or unsuccessful strategies.
Explain how these observations led to your conjecture.}

\section{Our Strategy, Conjecture, or Mathematical Claim}

\subsection{Precise statement}
\guidance{State the range of cases and all conditions under which your claim is
meant to hold. For a winning strategy, specify which player can use it, from
which positions, and whether it must work against every legal opponent move.
Repeated wins alone do not establish a guaranteed winning strategy.}

\textbf{Claim.} \placeholder{Write your conjecture or mathematical claim as a
complete, precise statement.}

\subsection{Strategy and intuition}
\placeholder{Describe the proposed strategy step by step, if applicable.
Explain the intuition behind it and how it follows from your observations.
If the claim is about counting, structure, or impossibility, explain the main
idea instead.}

\newpage
\section{Proof or Counterexample}
\guidance{Complete the appropriate route below and remove the unused route.
Experiments may suggest a claim; a proof must establish it for every case
covered by the statement. A valid counterexample must satisfy the assumptions
and contradict the claimed conclusion. Support your main claim with a proof
or disprove it with a counterexample.}

\subsection*{Route A: Proof of the claim}
\textbf{Proof.}
\placeholder{Give a clear, complete argument. State any facts or lemmas used
and justify the important steps. Use a suitable method, such as induction,
contradiction, a direct argument, an invariant, a counting argument, or a graph
argument. Choose the method that actually fits your claim.}

\placeholder{For induction, show the base case(s), state the induction
hypothesis, and prove the induction step. For a winning strategy, explain why
each prescribed move is legal, why every possible opponent response is
covered, and why the desired outcome is eventually reached. Remove this
prompt if it does not apply.}

\textbf{Conclusion of the proof.}
\placeholder{State exactly what has been established and for which cases.}

\subsection*{Route B: Counterexample to the conjecture}
\textbf{Original conjecture.}
\placeholder{Restate the conjecture being tested.}

\textbf{Counterexample.}
\placeholder{Give the specific starting state, construction, or sequence of
legal moves. Explain why it satisfies the assumptions and why the conclusion
fails. A loss caused by ignoring the proposed strategy is not a counterexample
to that strategy.}

\textbf{What the counterexample teaches us.}
\placeholder{Explain which part of the conjecture fails. If you propose a
revised claim, state it precisely and distinguish a proved result from a
conjecture that remains unproved.}

\subsection*{Final mathematical finding}
\placeholder{Summarise what you can now conclude about the game and whether
your main claim was proved or disproved. State the cases covered by your
result. Identify any further unanswered questions separately.}

\newpage
\section{Interaction with School Students and Feedback}
\guidance{The participants may be up to 15 years old. If the activity has not
yet taken place, label this section \textbf{Planned interaction}; describe the
plan and leave actual observations for later. Do not present planned outcomes
as observed results.}

\textbf{Introduction and guided play.}
\placeholder{Explain how you will introduce the rules simply, demonstrate a
round, and help students play. Give two or three guiding questions or hints
that encourage them to discover the mathematical idea.}

\textbf{Mathematics after playing.}
\placeholder{Write a brief, age-appropriate explanation of the idea you will
discuss after the students have explored the game.}

\textbf{Feedback and refinements, if available.}
\placeholder{Record what students or classmates found clear, confusing, or
interesting. Explain any changes made to the rules, materials, explanation,
or strategy. If no testing or interaction has occurred, state that clearly.}

\section{Contribution of Each Group Member}
\guidance{Use the same member order as on the first page. Describe each
person's specific contribution; shared tasks may appear in more than one row.
Examples include game design, construction, testing, mathematical analysis,
proof, illustrations, writing, and facilitating the activity.}

% longtable lets this table continue onto another page as entries grow.
\begin{longtable}{@{}P{1cm}P{4.3cm}P{\dimexpr\textwidth-5.3cm-4\tabcolsep\relax}@{}}
\toprule
\textbf{No.} & \textbf{Name / roll number} & \textbf{Specific contribution(s)} \\
\midrule
\endfirsthead
\toprule
\textbf{No.} & \textbf{Name / roll number} & \textbf{Specific contribution(s)} \\
\midrule
\endhead
\midrule
\multicolumn{3}{r}{\small Continued on the next page}\\
\endfoot
\bottomrule
\endlastfoot
1 & \placeholder{Member 1} & \placeholder{Contribution} \\
2 & \placeholder{Member 2} & \placeholder{Contribution} \\
3 & \placeholder{Member 3} & \placeholder{Contribution} \\
4 & \placeholder{Member 4} & \placeholder{Contribution} \\
5 & \placeholder{Member 5} & \placeholder{Contribution} \\
6 & \placeholder{Member 6} & \placeholder{Contribution} \\
7 & \placeholder{Member 7} & \placeholder{Contribution} \\
8 & \placeholder{Member 8} & \placeholder{Contribution} \\
9 & \placeholder{Member 9} & \placeholder{Contribution} \\
10 & \placeholder{Member 10} & \placeholder{Contribution} \\
\end{longtable}

\section{References and Acknowledgements}
\placeholder{List any books, websites, games, software, or other resources you
used. If you adapted an existing game or proof, credit the source and explain
your adaptation. If no external sources were used, state this.}

% Optional: add an appendix for extra trial data, code, or construction
% details. Keep the rules, mathematical question, claim, and complete
% proof/counterexample in the main report.

\end{document}
